\begin{center}
\begin{minipage}{0.98\linewidth}
\begin{tcolorbox}[colback=white,colframe=CStatement!25,boxrule=0.4pt,
  arc=1mm,left=3mm,right=3mm,top=3mm,bottom=3mm]
\centering
\textbf{中点代入为零，两端代入互为相反数}\par\medskip
\begin{tikzpicture}[font=\small,>=Stealth]

\begin{scope}[x=1.2cm,y=1.2cm]
\coordinate (P) at (0,3);\coordinate (H) at (1.44,1.08);\coordinate (Q) at (2.88,-.84);
\draw[CStatement,thick] (-.8,-.6)--(3.7,2.775) node[above left] {$f(x,y)=0$};
\draw[CExample,dashed] (-.8,2.4)--(1.1,3.825);
\draw[CTrap,dashed] (1.9,-1.575)--(3.7,-.225);
\draw[CProof,very thick] (P)--(Q);
\fill[CExample] (P) circle (2.3pt);\fill[CProof] (H) circle (2.3pt);\fill[CTrap] (Q) circle (2.3pt);
\node[above left,text=CExample] at (P) {$P:\ f(P)=D$};
\draw[gray] (H)--(2.0,1.0);\node[right,text=CProof] at (2.0,1.0) {$H:\ f(H)=0$};
\node[below right,text=CTrap] at (Q) {$P':\ f(P')=-D$};
\draw[CProof] (1.296,1.272)--(1.488,1.416)--(1.632,1.224);
\end{scope}
\end{tikzpicture}\par\medskip
$H=(P+P')/2$，所以 $f(H)=[f(P)+f(P')]/2=0$。图示 $D<0$，虚线分别经过两端点且平行于 $l$；$D$ 是代入值，不是实际距离。
\end{tcolorbox}
\end{minipage}
\end{center}
